\documentclass[11pt]{article}
\usepackage{amsmath,amssymb,amsthm,hyperref}
\newtheorem{theorem}{Theorem}\newtheorem{proposition}[theorem]{Proposition}
\newcommand{\E}{\mathcal E}
\title{Augmented polynomial pro-prefix semantics with backward position maps}
\date{16 November 2026}
\begin{document}\maketitle
\paragraph{Audit and scope.} The written repaired construction passed independent adversarial review and finite exact diagnostics within the stated scope. No novelty, compilation, or formal-verification claim is made. This adds an explicitly declared grade-map doctrine and content-discard structure. It constructs actual normalized native prefix towers for scope-atomic shared-frontier recurrence, and recovers their terminating part by the separately admitted first-hit clock construction. No infinite accumulated coefficient, arbitrary polynomial substitution, unrestricted pro-morphism preservation, or universal completion theorem is claimed. The strict time-indexed tower is the primary object; its associated pro-object is a second, weaker organizational interface.

\section{Changed doctrine and its recovery boundary}
Use the primitive processes and complete causal strategies of the retained source \emph{Causal Frontier Strategies}, and the countable row/coefficient assumptions of \emph{Countable Clock Trace and Causal Execution}. Context coefficients lie in $\E=\mathrm{Vect}_k/k$ with finite categorical products. Raw rows are countable and normalized. A primitive body is $f:A+U\to F_C(B+U)$, with finite arities, arbitrary small full shape families, and conditional frontier laws at every syntactic history, including zero histories. Each entire current frontier is one atomic sampling scope. Neither countability of the full shape family nor an almost-sure quotient of controls is assumed.

The additional content structure is a symmetric monoidal category $\mathcal D$ with a coherent augmentation $\epsilon_D:D\to I$ for every admitted coefficient, satisfying $\epsilon_I=1$, $\epsilon_{D\otimes E}=\epsilon_D\otimes\epsilon_E$ under the unitors, and $\epsilon_{D'}\delta=\epsilon_D$ for every admitted arrow $\delta:D\to D'$. Equivalently work in the symmetric monoidal category of augmented objects of an underlying symmetric monoidal category. Its unit is terminal among augmentation-preserving arrows. If the original content category was already affine this does not restrict its arrows; if augmentations have now been supplied, only the augmentation-compatible old subdoctrine is retained. Arbitrary old coefficient maps are not silently promoted to augmentation-preserving maps.

An enlarged polynomial grade is $G=(S,(P_s),(D_s))$, where $P_s$ is a finite ordered set and $D_s$ is augmented. A grade arrow $G\to G'$ has a full shape function $\alpha:S\to S'$, a backward position function
\[
 \beta_s:P'_{\alpha s}\longrightarrow P_s,
\]
and an actual augmented coefficient arrow $\delta_s:D_s\to D'_{\alpha s}$. No support restriction is made. The induced polynomial transformation sends $(s,x:P_s\to X)$ to $(\alpha s,x\beta_s)$. Composition uses $\beta_s\beta'_{\alpha s}$ and the composite coefficient arrows; identity uses identities. Thus these data form a category. The original reference is \emph{General Independent Interaction Criterion}, subsection ``Constant ordered polynomial content not generated by receipts'', retained source SHA-256 \texttt{508fffa3ed0cf355c9859f399677f4be3420d0dea877558163a8b0b2581d1c94}. It allows all component shape functions, including same-arity merges, but explicitly excludes binary-to-unary changes and arbitrary position reordering. Its coefficient-decorated form is the retained causal-frontier source. The old ordered arity-preserving grade arrows embed faithfully by the identity position bijection. New arity-changing and nonidentity-position arrows are genuinely added, so local fullness is not asserted.

For fixed $A,B$, let $\mathcal N_{A,B}$ be the category whose objects are normalized countable native context rows at each $a$, one full grade $G$, and bodies $W_a\to F_G(B)$. A map has the original exact row pushforward and actual context-coefficient arrows, together with one enlarged grade arrow, with the induced polynomial body square. These maps compose by their displayed data and preserve equality of distinct old native arrows. Varying $A,B$ along specified input/output functions can be treated with the same changing-endpoint squares; no arbitrary endofunction on lists replaces a pointwise output map.

\section{Full finite-prefix grades}
Define $P_0(U,C)$ to have the single seed shape, no variable positions, and coefficient $I$. The seed represents the as-yet-unexecuted external root; its input $a$ belongs to the process row, not to the grade. Recursively, at the polynomial level,
\[
 F_{P_{n+1}}(B)=F_C\bigl(B+U\times F_{P_n}(B)\bigr).       \tag{1}
\]
The coefficient of a shape is the tensor of all its primitive node coefficients, with $I$ for the seed. This recursive definition is a definition of all shapes and positions, not a restriction to process-realizable shapes.

Concretely, shapes are finite level-complete partial $C$-trees: each executed node carries a full $C$ shape; its ordered children are either settled variable $B$ positions or continuing edges labelled in $U$. Such an edge carries a further executed subtree unless it reaches the depth-$n$ frontier, where it is an active constant-labelled hole. If a tree has already terminated it has no active holes and appears unchanged at all larger horizons. There is no idle executed node after termination. The only empty-node tree is the seed at $n=0$. Nullary primitive nodes are genuine nonempty terminal trees at horizon one. Variable positions consist solely of the settled $B$ leaves, ordered by their planar addresses; active holes are constants in the shape, never variable positions. At a cut of depth zero the suffix is the original single-root computation with its external $A$-input retained in the raw state, not a $U$-labelled hole; the seed grade itself contains no external input value.

There are full grade maps
\[
 \tau_n:P_{n+1}\longrightarrow P_n\qquad(n\geq0).          \tag{2}
\]
For $n=0$, every shape maps to the seed, the backward position map has empty domain, and the coefficient arrow is its augmentation. For $n>0$, prune nodes of depth at least $n$ and replace their incoming $U$-labelled edges by active holes carrying those same labels. The backward map includes each retained settled-leaf address into the longer tree's settled positions. The coefficient arrow keeps every retained node factor and augments every removed factor, using the actual symmetric monoidal structural maps. Terminal trees of height at most $n$ are unchanged. These prescriptions are defined on every unused shape as well as on realized ones.

\begin{theorem}[Full graded truncation]
Equations (1)--(2) give an inverse tower of actual natural polynomial transformations with actual augmented coefficient maps. A composite truncation from height $m$ to $n$ is full tree pruning to depth $n$, the inclusion of retained settled positions, and augmentation of exactly the deleted node factors, independent of intermediate truncation levels. Full arity-preserving augmented collection maps $\alpha:C\to C'$ and state-label maps $u:U\to U'$ induce levelwise grade maps commuting with these truncations.
\end{theorem}
\begin{proof}
The shape map retains all existing labels on retained nodes and edges, including labels on edges newly exposed as holes. Every settled leaf in the pruned tree was already settled in the longer tree at the same address. This gives the displayed backward injection; composition of these inclusions is the inclusion of the final retained positions. For $n=0$ the empty function is unique. Consequently the grade map's induced polynomial transformation returns exactly the retained settled labels, naturally in $B$.

Choose the usual addressed order on node tensor factors. To delete a set of factors, use symmetry to group retained and deleted factors and tensor the identity on retained factors with the augmentations of the deleted factors. Symmetric monoidal coherence and monoidality of $\epsilon$ make this independent of the grouping and compatible with deletion in stages: every retained node receives its identity, and every deleted node its augmentation. This proves the asserted coefficient equality, including empty retained or deleted sets. It also proves augmentation preservation. Applying $\alpha$ does not change ordered node addresses, terminal tags, or arities, and $u$ merely changes constant edge labels. The augmented-arrow equation commutes the node maps past every deleted factor's augmentation. Thus both the position functions and actual coefficient arrows commute with pruning on the entire full shape families.
\end{proof}

\section{Normalized native prefix towers}
For a strategy $\sigma$ and external root $a$, let $H_{a,n}$ consist of all histories terminal by round $n$, together with all active histories at exact round $n$. A terminal history appears as itself at subsequent horizons, not as distinct strings of idle steps. At horizon zero there is one seed history. All syntactically possible histories are retained, including zero-probability ones. Weights are products of the actual successive frontier kernel probabilities. Context coefficients are finite products of all executed node coefficients in $\E$. Forgetting raw node witnesses gives a body in the full $P_n(U,C)$ with exactly the settled $B$ labels. Denote the resulting object of $\mathcal N_{A,B}$ by $X_n^\sigma$.

\begin{theorem}[Actual native recurrent tower]
Each $X_n^\sigma$ is normalized, and raw history truncation, actual context-product projection, and the full grade map $\tau_n$ define a native arrow
\[
 d_n:X_{n+1}^\sigma\longrightarrow X_n^\sigma.
\]
These form an inverse sequence; its composites agree with direct truncation through the explicitly identified canonical coefficient structural maps. All finite-horizon effects, coefficients, raw zero histories, active-hole labels and unused partial-tree grades are retained.
\end{theorem}
\begin{proof}
Normalization follows by induction. A terminal prefix has one unchanged successor with the same weight. An active prefix distributes its weight according to a normalized frontier kernel. Summing over countably many prefixes and draws gives total one. The raw map $d_n$ sends a terminal successor that was already terminal to itself; otherwise it removes exactly the last frontier draw. For a target active prefix $h$, its fibre is all its next frontier tuples, with total weight $m(h)\sum_{\mathbf w}q_h(\mathbf w)=m(h)$. For an already terminal prefix its fibre is the singleton persistent history. This proves the exact native pushforward equation, including every zero fibre, without dividing by prefix mass.

The context map is the actual product projection onto retained node factors, including the unique map to the empty product at horizon zero. The shape and settled labels of the truncated history are precisely those prescribed by $\tau_n$ and its backward position map, so the full body square commutes. The previous theorem proves actual content-coefficient compatibility and full-grade composition. Repeated context projections have identical node projections and hence are equal by product uniqueness. Together these verify all data and coherence of the inverse sequence.
\end{proof}

\begin{theorem}[Native comparison and cut compatibility]
For a source primitive native comparison with arity-preserving augmented collection map, state/output maps, and the local conditional-kernel equations $(r_h)_*q_h=q'_{Hh}$ at every source history, the induced maps $X_n^\sigma\to X_n^{\sigma'}$ are exact native maps and commute with every $d_n$. For a stopping rule bounded by $N$, they also commute with stopping-frontier cut/assembly at the retained marked grade on horizons $n\geq N$, whenever the stopping rules commute with $H$.
\end{theorem}
\begin{proof}
The raw history map preserves ordered addresses and the terminal/continuing tag and commutes with deletion of the last round. Induction on horizon proves exact pushforward: for each target next draw sum first over its source draw fibre, using the local kernel equation, and then over source prefixes, using the induction hypothesis. The same argument handles persistent terminal histories. No probability-zero history is omitted. Context maps are induced by the actual node maps into each target product, and commute with projection because both composites have the same target-node projections. The full grade square commutes by the full graded truncation theorem. These equations establish levelwise native naturality.

For a bound $N$ and horizons $n\geq N$, use the full marked partial-tree grade with shapes $(t,d)$ for every $t\in P_n$ and $0\leq d\leq\min(N,\operatorname{height}(t))$. The seed has height zero. The coefficient is the same node tensor and positions are the same settled leaves. Pruning from $n+1$ to $n$ leaves the marker $d$ unchanged: a shorter terminal tree is unchanged, and otherwise the pruned tree still has height at least $N\geq d$. Thus the full marked-grade map exists on unused shapes too. A bounded stopping cut retains the complete prefix and joint suffix, without separating draws inside an atomic scope. Their grafted finite prefix is exactly the original finite history. Splitting its kernel product and node product gives the source cut factorization; the marked partial-tree grade retains every permitted unused marked shape. Pruning before or after this grafting deletes exactly those nodes beyond the requested final horizon while retaining the same marker $d$ and prefix. The suffix horizon is $n-d$, and it is a joint forest over the entire retained frontier. The claim intentionally excludes pruning to a horizon earlier than the cut bound; a later stopping decision is not observable in such an earlier prefix. Content deletion commutes with grafting because both use the same retained node maps and augmentations of the same removed factors. Context deletion uses the same product projections. Thus the native squares commute, and repeated cuts/prunings are coherent by the same addresswise argument. No invertibility of the full forget-cut grade map or general scope-splitting law is inferred.
\end{proof}

\section{The pro-object and what its adequacy means}
An inverse sequence in $\mathcal N_{A,B}$ is in particular a pro-object. For clarity the pro-category here has the usual explicitly declared morphism sets
\[
 \operatorname{Hom}_{\mathrm{Pro}(\mathcal N)}((X_i),(Y_j))
 =\lim_j\operatorname*{colim}_i\operatorname{Hom}_{\mathcal N}(X_i,Y_j),
\]
with the transition maps induced by pre/postcomposition. Size is controlled by a fixed universe of small families. For fixed $A,B$, the preceding theorem constructs actual levelwise compatible maps, hence these pro-morphisms. With external maps $a:A\to A'$, $b:B\to B'$, first postcompose the source leaf labels by $b$ and pull back the target rows along $a$; the resulting towers and maps both lie in $\mathcal N_{A,B'}$. Constant objects embed fully faithfully, since for two constant diagrams the displayed formula reduces to the original hom set. No claim that this is the universal completion for causal feedback, or that all pro-morphisms are causal comparison maps, follows from that elementary calculation.

The strict indexed tower is retained as primary data. Passing only to a pro-isomorphism class permits reindexing and does not by itself retain a designated clock horizon. In particular extraction of an exact first-terminal layer is asserted natural for the displayed synchronous strategy comparisons, not for every arbitrary pro-morphism.

The passage from strict towers to pro-objects can also identify distinct native coefficient maps. For example, take deterministic always-continuing raw dynamics, with source root context coefficient $k\xrightarrow{1}k$ and every subsequent node coefficient $0\to k$. At horizon one the source coefficient is $k$; from horizon two onward it is $k\times_k0=0$. Let the target root coefficient be $k^2\xrightarrow{\pi_1}k$ with the same subsequent zero coefficients. The two root maps $x\mapsto(x,0)$ and $x\mapsto(x,x)$ are distinct augmented arrows, but their induced maps at every later source coefficient $0$ coincide. These zero vector coefficients may accompany weight-one raw witnesses: the scalar law and augmented-vector component are independent, and $0\to k$ is a valid augmented object. Take identical raw and full grade maps. These define distinct compatible strict-tower maps which yield the same pro-morphism, since they agree eventually. Thus no faithfulness of the strict-tower-to-pro passage is asserted; retaining the strict tower is essential when finite-stage native arrows are consumers.

\paragraph{Adequacy and its kernel.}
Every declared finite-horizon observation factors through its corresponding component $X_n$ simply by applying that observation to the retained native row; no finite-prefix data have been replaced by a limit coefficient. Equality of strict towers preserves all those data. At the scalar/raw level, the tower weights determine $q_h$ uniquely at positive histories by the ratio of successor and prefix weights. At a zero history every descendant prefix has zero weight, so the tower does not determine that history's contingent coupling table. Strategies differing only in such tables can have identical towers. The complete strategy therefore remains an explicit index if those controls are a consumer; no faithful reconstruction of strategies from the tower alone is claimed. Full native coefficient and grade maps remain visible in the tower regardless of those scalar zeros.

This supplies a recurrent semantics at the level of all finite native prefixes without forming an infinite tensor of executed content coefficients. For an always-continuing process each prefix is still normalized and records its growing partial tree, although the terminating process is empty. The object is not a single infinite-history weighted row. Existence of a scalar path-space measure or of a semantic limit in the original native category would be additional theorems, not consequences of the notation $\mathrm{Pro}$.

\section{Recovery of the clock-assembled terminating part}
For the clock-recovery theorem assume $A,B,U$ countable, as in the pinned countable-endpoint clock ambient. The prefix-tower theorems themselves require only small endpoint sets and countable raw histories reachable from each external input.

Let $T_n\subseteq P_n$ be the full subcollection of terminal trees of exact height $n\geq1$. It includes every such unused shape. Its positions are all its settled leaves, and its coefficient is the same finite node tensor. There is an arity-preserving full inclusion $T_n\to R_UC$, and the disjoint union of these full subcollections over $n$ is exactly $R_UC$. Every nonempty finite terminal tree has a unique positive height.

To form the $n$th clock row from $X_n$, apply the previously declared typed restriction to those raw histories whose shape lies in $T_n$, retype the resulting body to $T_n$, and then apply the grade inclusion into the entire $R_UC$. The restriction is not a native arrow from $X_n$; its effect mass can decrease. The latter grade extension is an actual native regrading at unchanged restricted raw row in the clock paper's subprobability native category, not in the normalized category $\mathcal N_{A,B}$ when the selected mass is less than one. All zero-weight terminal histories are retained, and no active prefix is assigned an invented terminal body.

\begin{theorem}[Recurrent/terminating assembly]
This procedure recovers exactly the first-hit clock compiler and its countable trace from \emph{Countable Clock Trace and Causal Execution}, at identity full final grade $R_UC$. Its resulting terminating process has the weights, raw trees and actual coefficients of causal branching execution. The recovery is natural for the displayed synchronous strategy comparisons and compatible with the bounded whole-frontier cuts of that theorem.
\end{theorem}
\begin{proof}
A raw prefix is in $T_n$ exactly when its first terminal round is $n$: terminal nodes are never advanced and breadth-first height counts executed rounds. Hence the selected raw set is the clock compiler's $T_{a,n}$, with identical finite products of transition weights and node coefficients. The full grade inclusion preserves every shape and coefficient in that layer and makes all other shapes of $R_UC$ unused but present. Tracing the countable clock collects each finite terminal history exactly once by its unique height. The raw bijection, actual product coefficient isomorphism and identity full-grade map are those of the clock theorem. Synchronous strategy comparisons preserve height and terminal tags and hence restrict to these strata; cut compatibility is its proved finite-round partition and countable assembly law. No arbitrary inverse-tower map is asserted to preserve a first-hit stratum.
\end{proof}

\paragraph{What has and has not been repaired.}
The original obstruction to actual content-prefix truncation is repaired by coherent augmentations and explicit backward position maps. The construction now has full unused graded prefixes, normalized exact native transition arrows, and a proved relationship to the terminating countable-trace semantics. It strictly enlarges the old grade-map interface: pruning a newly executed binary node can discard settled positions, which no arity-preserving old grade map can do. It does not preserve all native arrows of an arbitrary unaugmented original content category; does not remove the separately supplied typed restriction; and does not construct general unmarked shared polynomial substitution or an infinite accumulated effect. Those displaced obligations remain open rather than being silently solved in this changed doctrine.
\end{document}
